\documentclass[12pt]{article}
\setlength{\topmargin}{-0.5 in}
\setlength{\textheight}{9.5 in}
\setlength{\oddsidemargin}{-0.4 in}
\setlength{\textwidth}{7.0 in}
\usepackage{amsmath,graphicx,amssymb, diagbox, braket}
\usepackage{tikz}
\usepackage{braket}
\newcommand*\circled[1]{\tikz[baseline=(char.base)]{
		\node[shape=circle,draw,inner sep=2pt] (char) {#1};}}
\pagestyle{empty}
\begin{document}
	
\centerline{QC HW 6
	\hspace{0.4 in}V. Kamat
	\hspace{0.4 in}Spring 2026}
	\begin{quote}
		{\sf
		{\bf Directions.} Refer to the Week 6 notes (on quantum teleportation and the Deutsch--Jozsa circuit) for additional context surrounding each of the exercises below.  
		}
	\end{quote}
\begin{enumerate}
\item For each of the four cases $m_1m_2\in \{0,1\}^2$, show that Bob's qubit at the end of the quantum teleportation protocol is the same as Alice's qubit at the start. 
\item In the Deutsch--Jozsa algorithm, show that when $f$ is constant, the amplitude of any outcome $\ket{z} \neq |0\rangle^{\otimes n}$ is $0$. More precisely, prove that given $z\neq 0^n$
\[
\sum_{x\in\{0,1\}^n} (-1)^{x\cdot z} = 0.
\]
\item Consider a following application of the Deutsch-Jozsa circuit. Note that the state $\ket{\psi_3}$ attained during the execution of the Deutsch Jozsa circuit (ignoring, as before, the $\ket{-}$ in the final register) is independent of the definition of $f$. Suppose the ``promise'' of the problem is that $f$ is an \emph{affine} function given by $f(x)=a\cdot x\oplus b$, where $a,x\in \{0,1\}^n$, $a\cdot x$ is the $\text{mod}-2$ inner product, and $b\in \{0,1\}$. How would you determine $a$ using a single query to the black box $O_f$ implementing $f$? 
\end{enumerate}
\end{document}